\subsection{General non-negative measurable functions}

\begin{defn}
	let $f:E\gto[0,\infty]$ be meas.
	the \textbf{lebesgue integral} of $f$ is defined as:
	\begin{equation*}
		\int_{E}f \ = \ \sup\set{\int_{E}\phi:\phi \trm{ simple},0\leq\phi(x)\leq f(x) \ \fa x\in E}
	\end{equation*}
\end{defn}
by corr 13.5, if $f$ is simple, this agrees w our previous defn.

\begin{lm}
	if $0\leq f\leq g$, then:
	\begin{equation*}
		\int_{E}f \ \leq \ \int_{E}g
	\end{equation*}
	if $\lambda\in[0,\infty]$, then:
	\begin{equation*}
		\int_{E}\lambda f \ = \ \lambda\int_{E}f
	\end{equation*}
	again, 	we use the convention that $0\cdot \infty=0$.
\end{lm} \

\begin{pf}
	first part follows from defn bc if $0\leq\phi\leq f$, $\phi$ simple, then $\phi\leq g$.
	second follows for finite $\lambda$ since its true on simples.
	case $\lambda=\infty$ is an exercise.
\end{pf}

this is that other srs pf mentioned earlier.
\begin{thm}[title=Monotone Convergence Theorem]
	sps $f_{n}:E\gto[0,\infty]$ increasing mble fns.
	set $f$ st $f_{n}\sto f$ ptwise. then:
	\begin{equation*}
		\int_{E}f \ = \ \lim_{n\sto\infty}\int_{E}f_{n}
	\end{equation*}
\end{thm} \

\begin{pf}
	note the lim exists by lemma 14.2, and clearly:
	\begin{equation*}
		\lim_{n\sto\infty}\int_{E}f_{n} \ \leq \ \int_{E}f
	\end{equation*}
	now, let $0\leq\phi\leq f$ be simple, $\alpha\in(0,1)$. suffices to show:
	\begin{equation*}
		\lim_{n\sto\infty}\int_{E}f_{n} \ \geq \ \alpha\int_{E}\phi
	\end{equation*}
	since $\int_{E}f$ is the sup of RHS after taking $\alpha\sto1$.

	let $E_{k}=\set{x:f_{k}(x)\geq\alpha\phi(x)}$.
	since $f_{n}$ increasing, $\alpha\in(0,1)$, then
	$E_{k}\subseteq E_{k+1}$ and $\bigcup_{k}E_{k}=E$.
	let $\phi=\sum_{j=1}^{n}\beta_{j}\chi_{F_{j}}$. then:
	\begin{equation*}
		\lim_{n\sto\infty}\int_{E}f_{n}
		\ \geq \ \int_{E}f_{k} \ \geq \ \int_{E}f_{k}\chi_{E_{k}}
		\ \geq \ \int_{E}\alpha\phi\chi_{E_{k}}
		\ = \ \alpha\sum_{j=1}^{n}\beta_{j}m(F_{j}\cap E_{k})
	\end{equation*}
	this holds for all $k$, so as $k\sto\infty$, by continuity of measure:
	\begin{equation*}
		\lim_{n\sto\infty}\int_{E}f_{n}
		\ \geq \ \alpha\sum_{j=1}^{n}\beta_{j}m(F_{j})
		\ = \ \alpha\int_{E}\phi
	\end{equation*}
\end{pf} \

\begin{prop}
	we have additive linearity
\end{prop} \

\begin{pf}
	by lemma 11.3, since $f,g$ non-neg, we have
	$\phi_{n}\sto f$ and $\psi_{n}\sto g$ increasing seq of non-neg simples.
	thus:
	\begin{equation*}
		\int_{E}f+g \ = \ \int_{E}\lim_{n\sto\infty}\phi_{n}+\psi_{n}
		\ = \ \lim_{n\sto\infty}\int_{E}\phi_{n}+\psi_{n}
		\ = \ \lim_{n\sto\infty}\int_{E}\phi_{n}+\int_{E}\psi_{n}
		\ = \ \int_{E}f+\int_{E}g
	\end{equation*}
	by several invocations of MCT.
\end{pf} \

\begin{crll}
	if $f_{n}:E\gto[0,\infty]$, then:
	\begin{equation*}
		\int_{E}\sum_{n=1}^{\infty}f_{n} \ = \ \sum_{n=1}^{\infty}\int_{E}f_{n}
	\end{equation*}
\end{crll} \

\begin{lm}
	sps $f:E\gto[0,\infty]$ mble. then $\int_{E}f=0$ iff $f=0$ ae.
\end{lm} \

\begin{pf}
	if $f=0$ ae, then any simple $0\leq\phi\leq f$ is also 0 ae.
	indeed, if $\phi=\sum_{j=1}^{n}\alpha_{j}\chi_{E_{j}}$, then $\alpha_{j}>0 \implies m(E_{j})=0$.
	in particular, $\int_{E}\phi=0$, so $\int_{E}f=0$.

	now, sps $\int_{E}f=0$. let $E_{n}=\set{x:f(x)\geq1/n}$. then:
	\begin{equation*}
		0 \ = \ \int_{E}f \ \geq \ \int_{E}f\chi_{E_{n}}
		\ \geq \ \int_{E}\frac{\chi_{E_{n}}}{n}
		\ = \ \frac{m(E_{n})}{n} \ \geq \ 0
	\end{equation*}
	so $m(E_{n})=0$. now, let $A=\set{x:f(x)>0}$.
	clearly $A=\bigcup_{n}E_{n}$ and $E_{n}\subseteq E_{n+1}$, so:
	\begin{equation*}
		m(A) \ = \ \lim_{n\sto\infty}m(E_{n}) \ = \ 0
	\end{equation*}
	by continuity of measure.
\end{pf} \

\begin{thm}[title=Fatou's Lemma]
	let $f_{n}:E\sto[0,\infty]$ be mble. then:
	\begin{equation*}
		\int_{E}\liminf_{n\sto\infty}f_{n}
		\ \leq \ \liminf_{n\sto\infty}\int_{E}f_{n}
	\end{equation*}
\end{thm}

to rmbr which way the inequality goes, as an example,
take $E=\bR,f_{n}=\chi_{[n,n+1]}$. then:
\begin{equation*}
	\liminf_{n\sto\infty}f_{n}(x) \ = \ \lim_{n\sto\infty}f_{n}(x) \ = \ 0
	\ \implies \ \int_{\bR}\liminf_{n\sto\infty}f_{n} \ = \ 0
\end{equation*}
but:
\begin{equation*}
	\int_{\bR}f_{n} \ = \ m([n,n+1]) \ = \ 1
\end{equation*}
